\documentclass{beamer}
\usetheme{Madrid}
\usecolortheme{seagull}

% Packages
\usepackage[utf8]{inputenc}
\usepackage[vietnamese]{babel}
\usepackage{graphicx}
\usepackage{caption}
\usepackage{tikz}
\usepackage{amsmath}
\usepackage{booktabs}

% Font settings
\usefonttheme{professionalfonts}
\setbeamerfont{title}{size=\Large}
\setbeamerfont{frametitle}{size=\large}

% Footer với số slide
\setbeamertemplate{footline}[frame number]

% Thông tin chung
\title{Nghiên cứu BERT: Bidirectional Encoder Representations from Transformers}
\subtitle{Đồ án cuối kỳ - Hướng 2}
\author{Nhóm 13:\\
Lê Thái Ngọc - 23122012\\
Nguyễn Ngọc Khoa - 23122036\\
Nguyễn Lê Phúc Thắng - 22120332}
\institute{
Trường Đại học Khoa học Tự nhiên - ĐHQG HCM\\
Khoa Công nghệ Thông tin\\
Môn: Phương pháp toán cho Trí tuệ nhân tạo
}

\begin{document}

% Title slide
\begin{frame}
\titlepage
\end{frame}

% Outline
\begin{frame}{Nội dung trình bày}
\tableofcontents
\end{frame}

% Section 1: Giới thiệu
\section{Giới thiệu}
\begin{frame}{Bối cảnh và Động lực}
\begin{itemize}
    \item \textbf{Hạn chế của word embeddings tĩnh:}
    \begin{itemize}
        \item Word2Vec, GloVe: một từ → một vector cố định
        \item Không nắm bắt ngữ cảnh động (vd: "bank")
    \end{itemize}
    \item \textbf{Giải pháp trước BERT:}
    \begin{itemize}
        \item ELMo: Kết hợp 2 LSTM độc lập (shallow fusion)
        \item GPT: Chỉ học một chiều (left-to-right)
    \end{itemize}
    \item \textbf{Đột phá của BERT:}
    \begin{itemize}
        \item Học hai chiều sâu sắc (deeply bidirectional)
        \item Pre-train once, fine-tune for everything
    \end{itemize}
\end{itemize}
\end{frame}

\begin{frame}{Thành tích của BERT}
\begin{itemize}
    \item \textbf{GLUE Benchmark:}
    \begin{itemize}
        \item BERT\textsubscript{LARGE}: 80.5 điểm (+7.7 so với GPT)
        \item SOTA trên 11/11 tác vụ NLP
    \end{itemize}
    \item \textbf{SQuAD v1.1:}
    \begin{itemize}
        \item F1: 93.2 (vượt human performance 91.2)
        \item EM: 87.4
    \end{itemize}
    \item \textbf{Tác động:}
    \begin{itemize}
        \item Mở đầu kỷ nguyên Foundation Models
        \item Truyền cảm hứng cho GPT-3, T5, BART...
    \end{itemize}
\end{itemize}
\end{frame}

% Section 2: Kiến trúc
\section{Kiến trúc BERT}
\begin{frame}{Tổng quan Kiến trúc Transformer}
\begin{figure}
    \centering
    \includegraphics[width=0.7\textwidth]{figs/transformer_encoder.png}
    \caption{Kiến trúc Encoder trong Transformer (Vaswani et al., 2017)}
\end{figure}
\begin{itemize}
    \item BERT = Encoder-only Transformer
    \item Multi-Head Self-Attention + FFN
    \item Residual connections + Layer Norm
\end{itemize}
\end{frame}

\begin{frame}{Multi-Head Self-Attention}
\begin{columns}
\column{0.5\textwidth}
\begin{itemize}
    \item \textbf{Scaled Dot-Product:}
    $$\text{Attention}(Q,K,V) = \text{softmax}\left(\frac{QK^T}{\sqrt{d_k}}\right)V$$
    \item \textbf{Multi-Head:}
    \begin{itemize}
        \item h = 12 heads (BASE)
        \item Mỗi head học một góc nhìn
    \end{itemize}
\end{itemize}
\column{0.5\textwidth}
\begin{figure}
    \centering
    \includegraphics[width=\textwidth]{figs/multi_head_attention.png}
    \caption{Multi-Head Attention}
\end{figure}
\end{columns}
\end{frame}

\begin{frame}{Input Embeddings của BERT}
\begin{figure}
    \centering
    \includegraphics[width=0.8\textwidth]{figs/bert_input_embeddings.png}
    \caption{Ba loại embeddings được cộng để tạo input (Devlin et al., 2019)}
\end{figure}
\begin{itemize}
    \item \textbf{Token Embeddings:} WordPiece tokenization (30K vocab)
    \item \textbf{Segment Embeddings:} Phân biệt câu A và B
    \item \textbf{Position Embeddings:} Learned (không phải sinusoidal)
\end{itemize}
\end{frame}

\begin{frame}{So sánh BERT Base vs Large}
\begin{table}
\centering
\begin{tabular}{lcc}
\toprule
\textbf{Thông số} & \textbf{BERT\textsubscript{BASE}} & \textbf{BERT\textsubscript{LARGE}} \\
\midrule
Layers (L) & 12 & 24 \\
Hidden size (H) & 768 & 1024 \\
Attention heads (A) & 12 & 16 \\
Parameters & 110M & 340M \\
Training time & 4 days/16 TPUs & 4 days/64 TPUs \\
\bottomrule
\end{tabular}
\end{table}
\begin{itemize}
    \item Trade-off giữa hiệu năng và chi phí tính toán
    \item Cả hai đều vượt trội so với mô hình trước
\end{itemize}
\end{frame}

% Section 3: Pre-training
\section{Tác vụ Tiền huấn luyện}
\begin{frame}{Masked Language Model (MLM)}
\begin{figure}
    \centering
    \includegraphics[height=3.5cm]{figs/mlm_example.png}
    \caption{Ví dụ minh họa chiến lược masking 80-10-10}
\end{figure}
\begin{itemize}
    \item Che ngẫu nhiên 15\% tokens
    \item \textbf{Chiến lược 80-10-10:}
    \begin{itemize}
        \item 80\%: Thay bằng [MASK]
        \item 10\%: Thay bằng random token
        \item 10\%: Giữ nguyên
    \end{itemize}
    \item Giảm mismatch pre-training/fine-tuning
\end{itemize}
\end{frame}

\begin{frame}{Next Sentence Prediction (NSP)}
\begin{columns}
\column{0.6\textwidth}
\begin{itemize}
    \item \textbf{Mục tiêu:} Học quan hệ giữa câu
    \item \textbf{Dữ liệu:}
    \begin{itemize}
        \item 50\% IsNext (câu liên tiếp)
        \item 50\% NotNext (câu random)
    \end{itemize}
    \item \textbf{Dự đoán:} Từ [CLS] representation
    \item \textbf{Tranh cãi:} RoBERTa bỏ NSP → tốt hơn
\end{itemize}
\column{0.4\textwidth}
\begin{figure}
    \centering
    \includegraphics[width=\textwidth]{figs/nsp_example.png}
    \caption{Ví dụ NSP}
\end{figure}
\end{columns}
\end{frame}

% Section 4: Fine-tuning
\section{Fine-tuning và Thực nghiệm}
\begin{frame}{Quy trình Fine-tuning}
\begin{figure}
    \centering
    \includegraphics[width=0.8\textwidth]{figs/bert_finetuning.png}
    \caption{Fine-tuning BERT cho các tác vụ khác nhau}
\end{figure}
\begin{itemize}
    \item Giữ nguyên kiến trúc BERT
    \item Thêm task-specific output layer
    \item Train toàn bộ parameters end-to-end
\end{itemize}
\end{frame}

\begin{frame}{Thực nghiệm: Phân loại Cảm xúc IMDb}
\begin{itemize}
    \item \textbf{Dataset:} IMDb Movie Reviews
    \begin{itemize}
        \item Train: 4,000 reviews
        \item Validation: 1,000 reviews
    \end{itemize}
    \item \textbf{Baseline:} TF-IDF + Logistic Regression
    \item \textbf{Fine-tuning settings:}
    \begin{itemize}
        \item Learning rate: 2e-5
        \item Batch size: 16
        \item Epochs: 3
        \item Max length: 256 tokens
    \end{itemize}
\end{itemize}
\end{frame}

% Section 5: Kết quả
\section{Kết quả và Phân tích}
\begin{frame}{So sánh Hiệu năng}
\begin{table}
\centering
\begin{tabular}{lcccc}
\toprule
\textbf{Mô hình} & \textbf{Accuracy} & \textbf{Precision} & \textbf{Recall} & \textbf{F1} \\
\midrule
TF-IDF + LR & 0.855 & 0.86 & 0.86 & 0.85 \\
BERT (fine-tuned) & \textbf{0.894} & \textbf{0.89} & \textbf{0.90} & \textbf{0.89} \\
\midrule
\textbf{Cải thiện} & +3.9\% & +3.0\% & +4.0\% & +4.0\% \\
\bottomrule
\end{tabular}
\end{table}

\begin{figure}
    \centering
    \includegraphics[width=0.5\textwidth]{figs/accuracy_plot.png}
    \caption{Độ chính xác qua các epoch}
\end{figure}
\end{frame}

\begin{frame}{Ma trận Nhầm lẫn và Phân tích Lỗi}
\begin{columns}
\column{0.5\textwidth}
\begin{figure}
    \centering
    \includegraphics[width=\textwidth]{figs/confusion_matrix.png}
    \caption{Ma trận nhầm lẫn BERT}
\end{figure}
\column{0.5\textwidth}
\begin{itemize}
    \item \textbf{Balanced performance:}
    \begin{itemize}
        \item TP = TN = 447
        \item FP = FN = 53
    \end{itemize}
    \item \textbf{Khó khăn với:}
    \begin{itemize}
        \item Sarcasm/irony
        \item Negation phức tạp
        \item Domain-specific terms
    \end{itemize}
\end{itemize}
\end{columns}
\end{frame}

\begin{frame}{Trực quan hóa Không gian Embedding}
\begin{figure}
    \centering
    \includegraphics[width=0.8\textwidth]{figs/tsne_comparison.png}
    \caption{So sánh t-SNE: BERT (trái) vs TF-IDF (phải)}
\end{figure}
\begin{itemize}
    \item BERT: Phân tách rõ ràng positive/negative
    \item TF-IDF: Chồng chéo đáng kể
    \item → BERT học được biểu diễn ngữ nghĩa sâu sắc hơn
\end{itemize}
\end{frame}

\begin{frame}{Attention Visualization}
\begin{figure}
    \centering
    \includegraphics[width=0.7\textwidth]{figs/attention_heatmap.png}
    \caption{Heatmap attention cho câu "The movie was not bad"}
\end{figure}
\begin{itemize}
    \item Các head khác nhau focus vào patterns khác nhau
    \item Head 0: Quan hệ cú pháp
    \item Head 5: Negation ("not" → "bad")
\end{itemize}
\end{frame}

% Section 6: Kết luận
\section{Kết luận}
\begin{frame}{Tổng kết}
\begin{itemize}
    \item \textbf{Đóng góp chính của BERT:}
    \begin{itemize}
        \item Deeply bidirectional representations
        \item Pre-training objectives hiệu quả (MLM + NSP)
        \item Transfer learning paradigm cho NLP
    \end{itemize}
    \item \textbf{Ưu điểm:}
    \begin{itemize}
        \item SOTA trên nhiều benchmarks
        \item Fine-tuning đơn giản và hiệu quả
        \item Mở đường cho Foundation Models
    \end{itemize}
    \item \textbf{Hạn chế:}
    \begin{itemize}
        \item Chi phí tính toán cao
        \item Giới hạn 512 tokens
        \item Không tự nhiên cho text generation
    \end{itemize}
\end{itemize}
\end{frame}

\begin{frame}{Hướng Phát triển}
\begin{itemize}
    \item \textbf{Cải tiến hiệu quả:}
    \begin{itemize}
        \item RoBERTa: Training recipe tốt hơn
        \item ELECTRA: Replaced token detection
        \item DistilBERT: Knowledge distillation
    \end{itemize}
    \item \textbf{Xử lý văn bản dài:}
    \begin{itemize}
        \item Longformer: Sparse attention
        \item BigBird: Random + Window + Global
    \end{itemize}
    \item \textbf{Unified models:}
    \begin{itemize}
        \item T5: Text-to-Text framework
        \item GPT-3/4: Scaling to 175B+ params
    \end{itemize}
\end{itemize}
\end{frame}

% Thank you slide
\begin{frame}
\begin{center}
\Huge\textbf{Cảm ơn!}\\[1cm]
\Large Q\&A\\[1cm]
\normalsize
\textbf{Liên hệ:}\\
23122012@student.hcmus.edu.vn\\
23122036@student.hcmus.edu.vn\\
22120332@student.hcmus.edu.vn
\end{center}
\end{frame}

\end{document}